\chapter{相关理论与基础}

\section{路网表示学习}

路网的定量分析有着悠久的研究历史。Garrison 和 Kansky 等奠基性工作引入图论度量来描述交通网络的拓扑结构与连通模式\cite{Barthelemy2011}。经典方法采用回路数、连通性指标和可达性矩阵等图论特征来量化网络结构，并从静态图的角度研究城市形态对路网组织的影响。

Porta 等人\cite{PortaCrucitti2006a} 对城市街道网络进行了系统的双重图（Primal/Dual）分析，发现度分布、介数中心性等指标能有效区分不同城市形态的路网结构。Crucitti 等\cite{Crucitti2006} 进一步研究了路网中心性度量的分布规律，发现少数关键交叉口对全网连通性具有主导性作用，呈现出典型的幂律分布特征。这些研究奠定了从图拓扑角度理解路网的基础，但传统方法依赖手工设计的特征，缺乏对路网复杂语义信息的自动提取能力，难以满足现代智能交通系统对路网表示的需求。

随着深度学习的兴起，研究者开始探索从数据中自动学习路网嵌入表示的方法。早期工作将路网视为通用图，将网络嵌入技术（如 Node2Vec\cite{Grover2016}、DeepWalk\cite{Perozzi2014}）应用于路网，通过随机游走生成节点上下文并训练 Skip-gram 模型。RN2Vec\cite{wang2019rn2vec} 在此基础上引入了道路属性信息，以基于 token 的 skip-gram 方式学习路段嵌入，但这类方法仍未充分利用路网的图结构信息。

在时空交通预测方面，图神经网络与时序模型的结合开启了新的研究范式。Li 等\cite{li2018dcrnn} 将交通流建模为道路网络上的扩散过程，提出的 DCRNN 融合了双向随机游走图卷积与编码器-解码器架构。Yu 等\cite{yu2018stgcn} 提出时空图卷积网络（STGCN），首次以纯卷积结构替代循环神经网络处理时序依赖，显著提升了计算效率。Han 等\cite{han2024bigst} 提出线性复杂度时空图神经网络 BigST，实现了对十万量级大规模路网的高效交通预测，获得 VLDB 2024 最佳论文提名。

图神经网络（GNNs）的引入标志着路网表示学习的重要进展。Jepsen 等\cite{jepsen2019gcn_road} 将 GCN 应用于路网，通过图卷积聚合邻居路段的特征信息。GAT\cite{velickovic2018gat} 引入注意力机制后，在处理路段连接重要性差异方面展现出优势。关系融合网络（RFN）\cite{Jepsen2022} 通过结合原始图和对偶图的双视角来建模路网拓扑关系，在路段级别的分类任务上取得了显著提升。

Wu 等人提出的层次路网表示（HRNR）\cite{wu2020hrnr} 开创了显式建模路网多层次结构的先河。HRNR 构建了由功能区（Functional Zone）、结构区（Structural Zone）和路段（Road Segment）组成的三层神经架构，通过谱聚类初始化结构区划分，并利用轨迹数据监督功能区的学习。这一工作在多个下游任务上相比扁平 GNN 方法取得了显著提升，验证了层次结构建模的价值。

后续研究在 HRNR 的基础上取得了进一步进展。Toast\cite{chen2021toast} 提出了面向轨迹感知的路段预训练方法，通过轨迹数据设计多个自监督目标来增强路段嵌入。START\cite{chen2022start} 引入自监督轨迹感知表示框架，设计了更丰富的轨迹-路网交互学习目标。JCLRNT\cite{mao2022jclrnt} 提出路网-轨迹联合对比学习框架，通过多视角对比学习同时优化路网嵌入和轨迹嵌入。SARN\cite{li2023sarn} 引入自增强图对比学习，通过数据增强和对比损失提升路段表示的鲁棒性。HRoad\cite{lin2023hroad} 采用超图卷积来捕获路段间的高阶关系，进一步扩展了路网拓扑建模的边界。近年来，Zhou 等\cite{zhou2024garner} 提出基于第三地理学定律的路网表示学习框架（Garner），通过地理配置感知增强和谱负采样设计对比学习目标，使路段嵌入能够反映地理形态相似性规律，在多个路网基准任务上取得了显著提升。

上述方法虽然在路网表示学习领域取得了长足进步，但在层次建模能力上存在共同的局限。所有方法均在欧氏空间中运行，对层次关系的建模依赖于赋值矩阵或注意力机制，而非几何意义上的层次结构编码。HRNR 的三层架构虽然显式构建了层次结构，但层次关系通过软赋值矩阵以隐式方式表达，路段嵌入与其所属区域嵌入之间缺乏几何约束。其余方法（Toast、START、JCLRNT 等）虽引入了丰富的预训练目标，但其层次建模本质上仍是"扁平图上的特征增强"，而非"嵌入空间中的层次几何组织"。这一分析构成了本文引入双曲几何的核心动机。

\section{图神经网络}

图神经网络（Graph Neural Network, GNN）通过在图上执行消息传递（Message Passing）来学习节点表示，最早由 Scarselli 等\cite{scarselli2009gnn} 系统提出。Gilmer 等\cite{gilmer2017mpnn} 在此基础上提出了统一的消息传递神经网络（MPNN）框架，将主流 GNN 变体统一描述，推动了图学习理论的系统化发展。关于图神经网络的理论基础与最新进展，国内研究者亦开展了深入综述\cite{ma2022gnn_survey_cn,xu2020gcn_survey_cn,zhao2022large_gnn_cn}。对于节点 $v$ 的第 $l$ 层表示更新，一般形式为：
\begin{equation}
    \mathbf{h}_v^{(l)} = \text{UPDATE}^{(l)}\left(\mathbf{h}_v^{(l-1)}, \text{AGGREGATE}^{(l)}\left(\{\mathbf{h}_u^{(l-1)} : u \in \mathcal{N}(v)\}\right)\right)
\end{equation}
其中 $\mathcal{N}(v)$ 是节点 $v$ 的邻域，$\text{AGGREGATE}$ 和 $\text{UPDATE}$ 是可学习函数。不同 GNN 变体在聚合和更新函数的设计上有所差异。

GCN\cite{Kipf2017} 基于谱图理论，通过对图拉普拉斯矩阵的特征分解定义图上的卷积操作，其实际计算采用一阶切比雪夫多项式近似，形式为：
\begin{equation}
    \mathbf{H}^{(l+1)} = \sigma\left(\tilde{\mathbf{D}}^{-1/2}\tilde{\mathbf{A}}\tilde{\mathbf{D}}^{-1/2}\mathbf{H}^{(l)}\mathbf{W}^{(l)}\right)
\end{equation}
其中 $\tilde{\mathbf{A}} = \mathbf{A} + \mathbf{I}$ 为加自环的邻接矩阵，$\tilde{\mathbf{D}}$ 为对应度矩阵。GraphSAGE\cite{Hamilton2017} 提出了归纳式的邻居采样聚合方法，支持对大规模图上的新节点进行推断。GIN\cite{Xu2019} 从理论角度分析了 GNN 的表达能力，证明具有单射聚合函数的 GNN 与 Weisfeiler-Leman 图同构测试等价。

图注意力网络（GAT）\cite{Velickovic2018} 通过可学习的注意力系数对邻居的重要性进行加权，其注意力系数计算为：
\begin{equation}
    \alpha_{ij} = \frac{\exp(\text{LeakyReLU}(\mathbf{a}^T[\mathbf{W}\mathbf{h}_i \| \mathbf{W}\mathbf{h}_j]))}{\sum_{k \in \mathcal{N}(i)}\exp(\text{LeakyReLU}(\mathbf{a}^T[\mathbf{W}\mathbf{h}_i \| \mathbf{W}\mathbf{h}_k]))}
\end{equation}
其中 $\mathbf{a}$ 是注意力向量，$\|$ 表示拼接操作。GAT 在处理连接重要性差异显著的图时具有天然优势，这一特性使其特别适合路网分析——路网中不同道路类型之间的连接关系存在显著的重要性差异。

GNN 在交通领域的应用涵盖交通流预测、路段属性推断、轨迹建模等多个方向。HMIAN\cite{Sun2022} 通过层次映射与交互注意力机制融合多源时空数据用于交通预测。TrajRNE\cite{Scherer2023} 通过空间流卷积融合车辆轨迹信息来增强路段表示。VecCity\cite{zhang2025veccity} 提供了统一的路网表示学习基准框架，涵盖多城市数据集和标准化评估协议，为路网表示方法的公平比较提供了重要基础。邹等\cite{zou2025stgnn_survey_cn} 系统综述了基于图神经网络的复杂时空数据挖掘方法，深入梳理了交通预测、人流预测等城市计算任务的研究进展，Jin 等\cite{jin2024stgnn} 则从更宏观视角对城市计算中时空图神经网络的预测性学习进行了全面综述。

\section{双曲表示学习}

双曲几何是一类非欧几里得几何，在否定欧氏第五公设（平行公设）的基础上建立：在双曲空间中，过直线外一点可以引无数条与该直线不相交的直线。双曲空间具有恒定负曲率，其体积随半径呈指数增长，与树状层次结构中节点数量随深度指数增长的规律天然契合。

双曲空间有多种等价的数学模型。庞加莱球模型（Poincar\'e Ball Model）将双曲空间映射到开单位球内，满足共形性（角度保持），可视化直觉性强。Lorentz（超球面）模型将双曲空间嵌入 $(n+1)$ 维 Minkowski 时空的上双曲面上，数值稳定性更优，指数映射和对数映射具有封闭形式。两种模型之间可以通过解析变换相互转换，在不同应用场景中各有优势——本文第四章采用庞加莱球模型，第五章采用 Lorentz 模型，具体选择依据将在第三章理论分析中详细讨论。

Nickel 和 Kiela\cite{nickel2017poincare} 的开创性工作展示了庞加莱嵌入能够以显著低失真嵌入层次关系，将 WordNet 名词层次的嵌入维度从 GloVe 等欧氏方法的数百维降至 5 维。Ganea 等\cite{Ganea2018} 在庞加莱球模型中定义了双曲线性层、双曲循环神经网络等基本深度学习操作，开发了首个完整的双曲神经网络框架。Peng 等\cite{peng2022hyperbolic} 系统梳理了 Lorentz 模型的各类神经网络操作，证明了其数值稳定性优于庞加莱球模型，推动了 Lorentz 模型在实际应用中的广泛采用。

Chami 等\cite{Chami2019} 提出了双曲图卷积神经网络（HGCN），在超球面模型中推导图卷积操作，结合双曲注意力机制，在多个层次图数据集上链路预测误差最高降低 63.1\%。Liu 等\cite{Liu2019} 提出了双曲图神经网络（HGNN），在庞加莱球中实现了图卷积，在知识图谱层次学习任务上取得了优异表现。

Sala 等\cite{Sala2018} 从理论上分析了双曲嵌入与欧氏嵌入的权衡，证明双曲嵌入能以 $O(d)$ 维度以任意精度近似嵌入 $n$ 节点的树，而欧氏嵌入需要 $\Omega(\log n / d)$ 的失真下界。Khrulkov 等\cite{Khrulkov2020} 将双曲嵌入引入图像表示学习，在少样本学习和图像检索任务上验证了双曲空间对视觉特征层次性的建模能力，拓展了双曲表示学习的应用边界。

在双曲图神经网络的进一步发展中，Bachmann 等\cite{bachmann2020kgcn} 提出常曲率图卷积网络（$\kappa$-GCN），在统一的数学框架下实现了 GCN 向双曲空间、球面空间的推广，并支持乘积空间的组合。Chen 等\cite{chen2022fullyhyperbolic} 提出全双曲神经网络（FHNN），通过 Lorentz 旋转和提升变换将所有操作保持在双曲流形上，从根本上克服了现有方法过度依赖切空间线性化带来的精度损失。Yang 等\cite{yang2024hypformer} 提出 Hypformer——首个在双曲空间中完整实现的 Transformer 架构，借助 Lorentz 模型的线性自注意力机制将双曲 Transformer 扩展到十亿节点量级的大图。Sun 等\cite{sun2024lsenet} 提出 LSEnet，在 Lorentz 模型中引入可微分结构熵，实现了无需预设聚类数量的端到端双曲图聚类，进一步展示了 Lorentz 模型在图结构学习中的表达潜力。

蕴含锥（Entailment Cones）\cite{ganea2018entailment} 是在双曲空间中建模偏序关系的重要工具。给定双曲空间中的点 $\mathbf{q}$，其蕴含锥定义为以 $\mathbf{q}$ 为顶点的锥形区域，锥内的所有点被认为在语义上从属于 $\mathbf{q}$。锥的半孔径（Half-aperture）为：
\begin{equation}
    \omega(\mathbf{q}) = \sin^{-1}\left(\frac{2K}{\sqrt{\kappa}\|\tilde{\mathbf{q}}\|}\right)
\end{equation}
其中 $K$ 为常数，$\tilde{\mathbf{q}}$ 为 $\mathbf{q}$ 的空间分量。靠近原点的点（表示一般性概念）具有较宽的蕴含锥，远离原点的点（表示特定性概念）具有较窄的蕴含锥。

MERU\cite{desai2023meru} 将蕴含锥引入视觉-语言模型，在 Lorentz 模型中建立图像与文本之间的一般-特定关系。HyCoCLIP\cite{pal2025hycoclip} 在此基础上进一步提出了组合蕴含学习（Compositional Entailment Learning），通过蕴含锥对图像中整体-局部的组合层次关系进行建模，在层次分类任务上取得了最先进的性能，为本文第五章将该范式迁移至路网表示学习提供了核心灵感。

\section{对比学习与层次学习}


对比学习（Contrastive Learning）通过拉近正样本对的嵌入距离、推远负样本对来学习表示。Veličković 等\cite{velickovic2019dgi} 提出深度图信息最大化（DGI），通过最大化节点局部表示与全图摘要之间的互信息实现无监督图表示学习，开创了图对比学习的重要范式。You 等\cite{you2020graphcl} 系统研究了图数据增强策略对对比学习效果的影响，提出 GraphCL 框架并在多类分子图和社交网络数据上验证了其有效性。Zhu 等\cite{zhu2021gca} 进一步提出基于自适应增强的图对比学习（GCA），根据节点中心性和边重要性设计差异化的增强强度，缓解了统一增强策略对重要节点的破坏问题。

InfoNCE 损失\cite{oord2018infonce} 是最常用的对比学习目标，其形式为：
\begin{equation}
    \mathcal{L}_{\text{InfoNCE}} = -\log \frac{\exp(f(\mathbf{z}_i, \mathbf{z}_j)/\tau)}{\sum_{k} \exp(f(\mathbf{z}_i, \mathbf{z}_k)/\tau)}
\end{equation}
其中 $f(\cdot, \cdot)$ 为相似性函数，$\tau$ 为温度参数。在路网中，正对的定义可以基于拓扑邻接、功能相似性或轨迹共现等多种准则。


层次对比学习（Hierarchical Contrastive Learning）在标准对比学习的基础上，进一步考虑正负样本对的层次结构信息。在路网场景中，层次对比学习不仅需要拉近同层次的语义相似路段，还需要建立跨层次的关联——路段嵌入与其所属局部区域嵌入之间应保持较小的距离。

组合蕴含学习将对比损失与蕴含锥损失相结合，形成统一的层次学习框架：对比损失捕获同层次内的语义相似性，蕴含锥损失约束跨层次的包含关系。这一框架使模型能够同时学习层次结构的几何组织和层内的语义区分，为本文第五章将其迁移至路网提供了方法论基础。

\section{本章小结}

本章系统梳理了与本文研究密切相关的四个领域的代表性工作。通过上述综述可以看到：路网表示学习领域已从扁平 GNN 发展到层次建模，但仍局限于欧氏空间；双曲几何在层次结构嵌入方面具有显著的理论和实验优势；蕴含锥机制为显式建模偏序关系提供了有效工具；组合蕴含学习提供了将层次对比与蕴含约束统一的框架。这些研究背景共同构成了本文两项工作的理论与方法基础，也揭示了将双曲几何引入路网层次建模的研究空白——目前尚无工作系统性地利用双曲空间的几何特性来建模路网的层次结构。